\documentclass[article]{jss}

%% -- LaTeX packages and custom commands ---------------------------------------

%% recommended packages
\usepackage{thumbpdf,lmodern}

%% another package (only for this demo article)
\usepackage{framed}

%% Packages from FCVAR.m
\usepackage{amsmath}
\usepackage{amsfonts}
\usepackage{float}   % used for better positioning of figures
\usepackage{mathrsfs} % for script fonts
% \usepackage[FIGTOPCAP]{subfigure}     %FIGTOPCAP puts captions on top of subfigures
\usepackage{subfigure}     %FIGTOPCAP puts captions on top of subfigures (not for JSS)
% \usepackage{subfig}

%% new custom commands
\newcommand{\class}[1]{`\code{#1}'}
\newcommand{\fct}[1]{\code{#1()}}

%% new custom commands from FCVAR.m
\def\lrt{{\rm LR}_T}


%% -- Article metainformation (author, title, ...) -----------------------------

%% - \author{} with primary affiliation
%% - \Plainauthor{} without affiliations
%% - Separate authors by \And or \AND (in \author) or by comma (in \Plainauthor).
%% - \AND starts a new line, \And does not.
\author{Lealand Morin\\University of Central Florida
   \And Harry J. Paarsch\\University of Central Florida}
\Plainauthor{Lealand Morin, Harry J. Paarsch}

%% - \title{} in title case
%% - \Plaintitle{} without LaTeX markup (if any)
%% - \Shorttitle{} with LaTeX markup (if any), used as running title
\title{Discretizing the Distribution of a Population following a Discrete-time Markov Chain on a Continuous State Space}
\Plaintitle{Discretizing the Distribution of a Population following a Discrete-time Markov Chain on a Continuous State Space}
\Shorttitle{Discretizing Continuous-state Markov Chains}

%% - \Abstract{} almost as usual
%% I figure that it is best to explain what this package does in plain English,  
%% in case any normal people get past the title and want to know what this package does. 
\Abstract{
  Package \pkg{discretizeCtsDTMC} creates a discrete approximation to a Markov
  process defined on a continuous state space in discrete time.
  Once the state space is discretized, \pkg{discretizeCtsDTMC} provides tools
  to estimate the transition matrices and analyze the Markov process.
  It is used to model a population of individuals, each following a  
  continuous-state Markov process in discrete time. \\

  We'll choose a better name later.


}

%% - \Keywords{} with LaTeX markup, at least one required
%% - \Plainkeywords{} without LaTeX markup (if necessary)
%% - Should be comma-separated and in sentence case.
\Keywords{Markov chain, Markov process, discretization}
\Plainkeywords{Markov chain, Markov process, discretization}

%% - \Address{} of at least one author
%% - May contain multiple affiliations for each author
%%   (in extra lines, separated by \emph{and}\\).
%% - May contain multiple authors for the same affiliation
%%   (in the same first line, separated by comma).
\Address{
  Lealand Morin\\
  Department of Economics\\
  University of Central Florida\\
  E-mail: \email{Lealand.Morin@ucf.edu}
  \phantom{0}\\
  Harry J. Paarsch\\
  Department of Economics\\
  University of Central Florida\\
  E-mail: \email{Harry.Paarsch@ucf.edu}
}

\begin{document}


%% -- Introduction -------------------------------------------------------------

%% - In principle "as usual".
%% - But should typically have some discussion of both _software_ and _methods_.
%% - Use \proglang{}, \pkg{}, and \code{} markup throughout the manuscript.
%% - If such markup is in (sub)section titles, a plain text version has to be
%%   added as well.
%% - All software mentioned should be properly \cite-d.
%% - All abbreviations should be introduced.
%% - Unless the expansions of abbreviations are proper names (like "Journal
%%   of Statistical Software" above) they should be in sentence case (like
%%   "generalized linear models" below).

\section[Markov Chains in R]{Markov Chains in \proglang{R}} \label{sec:intro}

% JSS Notes:
%\begin{leftbar}
%The introduction is in principle ``as usual''. However, it should usually embed
%both the implemented \emph{methods} and the \emph{software} into the respective
%relevant literature. For the latter both competing and complementary software
%should be discussed (within the same software environment and beyond), bringing
%out relative (dis)advantages. All software mentioned should be properly
%\verb|\cite{}|d. 
% 
%For writing about software JSS requires authors to use the markup
%\verb|\proglang{}| (programming languages and large programmable systems),
%\verb|\pkg{}| (software packages), \verb|\code{}| (functions, commands,
%arguments, etc.). If there is such markup in (sub)section titles (as above), a
%plain text version has to be provided in the {\LaTeX} command as well. Below we
%also illustrate how abbrevations should be introduced and citation commands can
%be employed. See the {\LaTeX} code for more details.
% \end{leftbar}

This article illustrates how to creates a discrete approximation to a Markov
  process defined on a continuous state space in discrete time.
  Once the state space is discretized, \pkg{discretizeCtsDTMC} provides tools
  to estimate the transition matrices and analyze the Markov process.
  It is used to model a population of individuals, each following a  
  continuous-state Markov process in discrete time.
% 


% 
The R packages \citep{R} available for working with Markov chains is as follows...

\begin{figure}[H]
  \centering
  \includegraphics[scale = .6, keepaspectratio=true]{Figures/Rlogo.png}
  \caption{Caption goes here}
  \label{fig:Rlogo}
\end{figure}



% JSS Style Notes:
%% -- Manuscript ---------------------------------------------------------------
%% - In principle "as usual" again.
%% - When using equations (e.g., {equation}, {eqnarray}, {align}, etc.
%%   avoid empty lines before and after the equation (which would signal a new
%%   paragraph.
%% - When describing longer chunks of code that are _not_ meant for execution
%%   (e.g., a function synopsis or list of arguments), the environment {Code}
%%   is recommended. Alternatively, a plain {verbatim} can also be used.
%%   (For executed code see the next section.)

\section{Model} \label{sec:model}

% JSS notes:
%\begin{leftbar}
%Note that around the \verb|{equation}| above there should be no spaces (avoided
%in the {\LaTeX} code by \verb|%| lines) so that ``normal'' spacing is used and
%not a new paragraph started.
%\end{leftbar}


Harry, put all that good stuff here. 


% JSS Notes:
%% -- Illustrations ------------------------------------------------------------
%%
%% - Virtually all JSS manuscripts list source code along with the generated
%%   output. The style files provide dedicated environments for this.
%% - In R, the environments {Sinput} and {Soutput} - as produced by Sweave() or
%%   or knitr using the render_sweave() hook - are used (without the need to
%%   load Sweave.sty).
%% - Equivalently, {CodeInput} and {CodeOutput} can be used.
%% - The code input should use "the usual" command prompt in the respective
%%   software system.
%% - For R code, the prompt "R> " should be used with "+  " as the
%%   continuation prompt.
%% - Comments within the code chunks should be avoided - these should be made
%%   within the regular LaTeX text.

%But first... Examples of displaying code in JSS style. 
%
%The data can be loaded by
%%
%\begin{CodeChunk}
%\begin{CodeInput}
%R> data("quine", package = "MASS")
%\end{CodeInput}
%\end{CodeChunk}
%%
%and a basic frequency distribution of the response variable is displayed in
%Figure~\ref{fig:quine}.

%\begin{leftbar}
%For code input and output, the style files provide dedicated environments.
%Either the ``agnostic'' \verb|{CodeInput}| and \verb|{CodeOutput}| can be used
%or, equivalently, the environments \verb|{Sinput}| and \verb|{Soutput}| as
%produced by \fct{Sweave} or \pkg{knitr} when using the \code{render_sweave()}
%hook. Please make sure that all code is properly spaced, e.g., using
%\code{y = a + b * x} and \emph{not} \code{y=a+b*x}. Moreover, code input should
%use ``the usual'' command prompt in the respective software system. For
%\proglang{R} code, the prompt \code{"R> "} should be used with \code{"+  "} as
%the continuation prompt. Generally, comments within the code chunks should be
%avoided -- and made in the regular {\LaTeX} text instead. Finally, empty lines
%before and after code input/output should be avoided (see above).
%\end{leftbar}

% JSS format for figures:
%A basic frequency distribution of the response variable is displayed in
%%Figure~\ref{fig:quine}.
%
%\begin{figure}[t!]
%\centering
%\includegraphics{delete-this-image}
%\caption{\label{fig:quine} Frequency distribution for number of days absent
%from school.}
%\end{figure}


%As a first model for the \code{quine} data, we fit the basic Poisson regression
%model. (Note that JSS prefers when the second line of code is indented by two
%spaces.)
%%
%\begin{CodeChunk}
%\begin{CodeInput}
%R> m_pois <- glm(Days ~ (Eth + Sex + Age + Lrn)^2, data = quine,
%+    family = poisson)
%\end{CodeInput}
%\end{CodeChunk}
%%
%Hence, the full summary of that model is shown below.
%%
%\begin{CodeChunk}
%\begin{CodeInput}
%R> summary(m_nbin)
%\end{CodeInput}
%\begin{CodeOutput}
%Call:
%glm.nb(formula = Days ~ (Eth + Sex + Age + Lrn)^2, data = quine, 
%    init.theta = 1.60364105, link = log)
%
%Coefficients: (1 not defined because of singularities)
%            Estimate Std. Error z value Pr(>|z|)    
%(Intercept)  3.00155    0.33709   8.904  < 2e-16 ***
% ...
%\end{CodeOutput}
%\end{CodeChunk}
%
%
%
%Here is another example of code:
%\begin{Code}
%glm(formula, data, subset, na.action, weights, offset,
%  family = gaussian, start = NULL, control = glm.control(...),
%  model = TRUE, y = TRUE, x = FALSE, ...)
%\end{Code}


%\begin{leftbar}
%As the synopsis above is a code listing that is not meant to be executed,
%one can use either the dedicated \verb|{Code}| environment or a simple
%\verb|{verbatim}| environment for this. Again, spaces before and after should be
%avoided.
%\end{leftbar}


%% -- Main Program ---------------------------------------------------------
% \section{Main Program} \label{sec:main}
% Renamed as Example session

\section{Example} \label{sec:example}

A demonstration of analysis is shown in \verb|discCtsDTMC_demo.R| and it serves as an example of what a typical session of model specification, estimation and testing can include. This procedure includes the following steps:

\begin{enumerate}
\item Organizing data
\item Choosing estimation options
\item Lag selection
\item Model estimation
\item Hypothesis testing
\end{enumerate}


\subsection{Organizing data} \label{subsec:data}


\subsection{Choosing options} \label{subsec:options}


\subsection{Lag-order selection}  \label{subsec:lags}


\subsection{Model estimation}  \label{subsec:estimation}


\subsection{Hypothesis testing} \label{subsec:testing}




%\begin{leftbar}
%Finally, there might be a reference to a \verb|{table}| such as
%Table~\ref{tab:overview}. Usually, these are placed at the top of the page
%(\verb|[t!]|), centered (\verb|\centering|), with a caption below the table,
%column headers and captions in sentence style, and if possible avoiding vertical
%lines.
%\end{leftbar}

%\begin{table}[t!]
%\centering
%\begin{tabular}{lllp{7.4cm}}
%\hline
%Type           & Distribution & Method   & Description \\ \hline
%GLM            & Poisson      & ML       & Poisson regression: classical GLM,
%                                           estimated by maximum likelihood (ML) \\
%Zero-augmented & Poisson      & ML       & Zero-inflated Poisson (ZIP),
%                                           hurdle Poisson \\
%               & NB           & ML       & Zero-inflated NB (ZINB),
%                                           hurdle NB \\ \hline
%\end{tabular}
%\caption{\label{tab:overview} Overview of various count regression models. The
%table is usually placed at the top of the page (\texttt{[t!]}), centered
%(\texttt{centering}), has a caption below the table, column headers and captions
%are in sentence style, and if possible vertical lines should be avoided.}
%\end{table}




%% -- Summary/conclusions/discussion -------------------------------------------

\section{Summary and discussion} \label{sec:summary}

This is a good package because...


%% -- Optional special unnumbered sections -------------------------------------

\section*{Computational details}

%\begin{leftbar}
%If necessary or useful, information about certain computational details
%such as version numbers, operating systems, or compilers could be included
%in an unnumbered section. Also, auxiliary packages (say, for visualizations,
%maps, tables, \dots) that are not cited in the main text can be credited here.
%\end{leftbar}

The results in this paper were obtained using
\proglang{R}~3.5.1. 
with the \pkg{discreteCtsDTMC} package Version 0.0.0.9000. 
\proglang{R} itself
and all packages used are available from the Comprehensive
\proglang{R} Archive Network (CRAN) at
\url{https://CRAN.R-project.org/}.


The development version of this package is available by using the \pkg{devtools} package, with which the latest version can be installed by 

\verb|devtools::install_github(LeeMorinUCF/discreteCtsDTMC)|. 



\section*{Acknowledgments}

%\begin{leftbar}
%All acknowledgments (note the AE spelling) should be collected in this
%unnumbered section before the references. It may contain the usual information
%about funding and feedback from colleagues/reviewers/etc. Furthermore,
%information such as relative contributions of the authors may be added here
%(if any).
%\end{leftbar}

To be amended: We are particularly grateful to .... For useful comments and helpful suggestions
on earlier versions of this paper, we thank ....



%% -- Bibliography -------------------------------------------------------------
%% - References need to be provided in a .bib BibTeX database.
%% - All references should be made with \cite, \citet, \citep, \citealp etc.
%%   (and never hard-coded). See the FAQ for details.
%% - JSS-specific markup (\proglang, \pkg, \code) should be used in the .bib.
%% - Titles in the .bib should be in title case.
%% - DOIs should be included where available.

\bibliography{references}

%\begin{leftbar}
%Note that References have to be written in Title Case, as opposed to sentence case, 
%sometimes contrary to the convention in the journals themselves. 
%\end{leftbar}


%% -- Appendix (if any) --------------------------------------------------------
%% - After the bibliography with page break.
%% - With proper section titles and _not_ just "Appendix".

% \newpage

% \begin{appendix}

% \section{More technical details} \label{app:technical}

%\begin{leftbar}
%Appendices can be included after the bibliography (with a page break). Each
%section within the appendix should have a proper section title (rather than
%just \emph{Appendix}).
%
%For more technical style details, please check out JSS's style FAQ at
%\url{https://www.jstatsoft.org/pages/view/style#frequently-asked-questions}
%which includes the following topics:
%\begin{itemize}
%  \item Title vs.\ sentence case.
%  \item Graphics formatting.
%  \item Naming conventions.
%  \item Turning JSS manuscripts into \proglang{R} package vignettes.
%  \item Trouble shooting.
%  \item Many other potentially helpful details\dots
%\end{itemize}
%\end{leftbar}

% Technical details go here. 


%\section[Using BibTeX]{Using \textsc{Bib}{\TeX}} \label{app:bibtex}
%
%\begin{leftbar}
%References need to be provided in a \textsc{Bib}{\TeX} file (\code{.bib}). All
%references should be made with \verb|\cite|, \verb|\citet|, \verb|\citep|,
%\verb|\citealp| etc.\ (and never hard-coded). This commands yield different
%formats of author-year citations and allow to include additional details (e.g.,
%pages, chapters, \dots) in brackets. In case you are not familiar with these
%commands see the JSS style FAQ for details.
%
%Cleaning up \textsc{Bib}{\TeX} files is a somewhat tedious task -- especially
%when acquiring the entries automatically from mixed online sources. However,
%it is important that informations are complete and presented in a consistent
%style to avoid confusions. JSS requires the following format.
%\begin{itemize}
%  \item JSS-specific markup (\verb|\proglang|, \verb|\pkg|, \verb|\code|) should
%    be used in the references.
%  \item Titles should be in title case.
%  \item Journal titles should not be abbreviated and in title case.
%  \item DOIs should be included where available.
%  \item Software should be properly cited as well. For \proglang{R} packages
%    \code{citation("pkgname")} typically provides a good starting point.
%\end{itemize}
%\end{leftbar}

% \end{appendix}

%% -----------------------------------------------------------------------------


\end{document}
